\subsection{A literal finite braid bridge for the branch-order comparison}

\begin{result}{A.4.26}{Proved}{One explicit pure gauge identifies all three braid-generator images}
Let $a,b,d$ be the three standard chain twists on the double
cover of a disc branched at four ordered points.  This is a genus-one
surface with two boundary circles.  Write
\[
 V=(aba)^2=(ab)^3,\qquad e=V^2,\qquad
 \omega=(abd)^4,
 \qquad g=a^{8c}V^{-4c},\quad c=c_2(F).                  \tag{ABR.1}
\]
For the quilt action of A.6.1 choose the pants curves $(a,e)$,
with the two independent adaptations specified in the proof.
For every $F=(\lambda,f)\in\widehat{\mathrm{GT}}$, its three
images are
\[
\begin{aligned}
 \Phi(a)&=a^\lambda,\\
 \Phi(b)&=a^{-8c}f(b^2,a^2)b^\lambda
                         f(b^2,a^2)^{-1}a^{8c},\\
 \Phi(d)&=f(d,V^2)d^\lambda f(d,V^2)^{-1}.
\end{aligned}                                                    \tag{ABR.2}
\]
The single explicit conjugation $\Psi=\operatorname{Inn}(g)\Phi$
has images
\[
\begin{aligned}
 \Psi(a)&=a^\lambda,\\
 \Psi(b)&=f(a^2,b^2)^{-1}b^\lambda f(a^2,b^2),\\
 \Psi(d)&=f(y_3,d^2)^{-1}d^\lambda f(y_3,d^2),
 \qquad y_3=ba^2b=a^{-2}V.
\end{aligned}                                                    \tag{ABR.3}
\]
These are precisely the three genuine braid formulas on $B_4$.
The gauge $g$ belongs to the pure subgroup.  Consequently this
comparison also gives the actual branch-order diagram
$P_4/\langle\omega\rangle\to B_4/\langle\omega\rangle$
with all four branch labels retained.

The comparison is a finite list of proword evaluations and braid
identities.  For any finite image $q:B_4\to Q$, its complete
additional coefficient test, beyond the local weights themselves,
is the identity
\[
\begin{split}
 f(q(d),q(V)^2)
   =q(\omega)^{-2c}q(V)^{4c}
       f(q(d)^2,q(V))q(d)^{4c}q(a)^{-4c}.
\end{split}                                                     \tag{ABR.4}
\]
This is the $(d,V)$ evaluation of (TA.4), with its actual
mixed-braid homomorphism specified below.  All gauges and
equalities in (ABR.2)--(ABR.3) are then the images under $q$ of
the displayed words.  Thus this finite comparison is supplied
by explicit morphisms; it does not use a classification of
outer automorphisms of a finite mapping-class image.
\end{result}

\begin{proof}
\par\medskip\noindent\textbf{The exact surface and the two adaptations.}\quad
Place the four branch points in order on a diameter of a disc.
The lifts of the three adjacent half twists are the three positive
chain twists $a,b,d$.  Their curves satisfy
$i(a,b)=i(b,d)=1$ and $i(a,d)=0$.  The inverse image of a small
disc surrounding the first three branch points is a one-holed
torus.  Its boundary is the separating curve denoted $e$.
The full twist of these three branch points lifts to $V$;
on the boundary annulus the double cover has degree two, so its
square is the positive boundary twist.  Equivalently the literal
two-chain relation gives $D_e=(ab)^6=V^2$.
We use the same letter $e$ for this positive twist in (ABR.1).

The curves $a,e$ cut the two-boundary surface into two pairs of
pants.  Inside the one-holed torus, replacement of $a$ by $b$
is the one-holed-torus S-move.  Cutting instead along $a$ gives
a four-holed sphere; the curves $e,d$ there intersect twice,
and replacement of $e$ by $d$ is an A-move.  This describes
both moves on the same marked surface and the same chain,
not merely their abstract local types.

Start with any quilt on the pants decomposition $(a,e)$.
Its half-adjustment lattice has separate coordinates at $a$
and $e$.  Altering the $a$ coordinate makes the S-move
$a\to b$ adapted, with adjustment integer zero.  It does not
change the A-move adjustment: $a$ is retained by that move,
and its collar is disjoint from both replacement curves $e,d$.
Altering the $e$ coordinate then makes $e\to d$ adapted.  It
does not change the S-move adjustment, because $e$ is a retained
boundary of its one-holed-torus support and is disjoint from
both $a,b$.  The two collar adjustments commute, since their
supports are disjoint.  These statements are also exactly the
retained-curve adjustment rule used in (TA.9), with the
collar construction of \cite[\S7, Definition~7.2 and
Lemmas~7.1--7.3]{ET_NS2000}. Hence one quilt
has both required adjustment integers zero.

The cuff $a$ has image $a^\lambda$.  The adapted S-edge has
weight
\[
 a^{-8c}f(b^2,a^2)b^{8c}(aba)^{2\nu},
 \qquad \nu=(\lambda-1)/2.
\]
The two final factors commute with $b$: the first is a power of
$b$, and $(aba)^2=V$ is central in $\langle a,b\rangle$.
Its conjugation of $b^\lambda$ therefore gives the middle
formula in (ABR.2).  The adapted A-edge replacing $e$ by $d$
has weight $f(d,e)$, so its conjugation formula gives the last
line.  This proves (ABR.2), with both geometric adjustment
choices recorded.

\par\medskip\noindent\textbf{The one mixed-braid substitution.}\quad
In $B_4$ the exact discrete identities are
\[
 [a,V]=[b,V]=[a,d]=1,
 \qquad (dV)^2=\omega a^2,
 \qquad \omega\in Z(B_4).                               \tag{ABR.5}
\]
The mixed subgroup $\langle u,w^2\rangle\leq B_3$
has the presentation
\[
 \langle u,v\mid(uv)^2=(vu)^2\rangle,
 \qquad v=w^2.                                         \tag{ABR.6}
\]
It is the inverse image of the stabilizer of one strand;
use the right-coset representatives $1,w,wu$.  Their four
nontrivial Schreier generators can be denoted $u,v,k,t$.
Rewriting $uwu=wuw$ at the three representatives gives
$t=u$, $k=u^{-1}vu$, and $kvu=t kv$.
After these substitutions the last equality is
$(vu)^2=(uv)^2$.  This proves the presentation (ABR.6),
including that it has no additional relations.
The map $u\mapsto d$, $v\mapsto V$ is well-defined, since
$(dV)^2=\omega a^2$ commutes with both $d$ and $V$.
The $(p,q)=(d,V)$ row of (TA.3)--(TA.4) is precisely this
homomorphism, not a choice of a square root of $V$ in $B_4$.
Applying that row yields
\[
 f(d,V^2)=\omega^{-2c}V^{4c}
                f(d^2,V)d^{4c}a^{-4c}.                 \tag{ABR.7}
\]
The rightmost powers of $d,a$ commute with $d^\lambda$,
and the power of $\omega$ is central.  Thus (ABR.2) becomes
\[
 \Phi(d)=V^{4c}f(d^2,V)d^\lambda
                     f(d^2,V)^{-1}V^{-4c}.             \tag{ABR.8}
\]

Now $g$ commutes with $a$, so its first conjugated image is
$a^\lambda$.  On the second image,
\[
\begin{split}
 g\,a^{-8c}f(b^2,a^2)b^\lambda
                f(b^2,a^2)^{-1}a^{8c}\,g^{-1}
 =f(b^2,a^2)b^\lambda f(b^2,a^2)^{-1},
\end{split}                                                     \tag{ABR.9}
\]
because $V$ commutes with every element of the $a,b$ subgroup.
On the third image, (ABR.8) gives
\[
\begin{split}
 g\,V^{4c}f(d^2,V)d^\lambda
                  f(d^2,V)^{-1}V^{-4c}\,g^{-1}
 =f(d^2,V)d^\lambda f(d^2,V)^{-1}.
\end{split}                                                     \tag{ABR.10}
\]
Here the two powers of $V$ cancel adjacently, and $a$ commutes
with both arguments $d^2,V$ and with $d$.
These calculations prove all three conjugations, with one gauge.

For the exact change from $V$ to the third braid word, the
faithful Artin identities give $V=a^2y_3$, where
$y_3=ba^2b$.  The element $a^{-2}$ commutes with $V$ and $d^2$.
Consider the actual map
$\widehat F(X,Y)\times\widehat{\mathbb Z}(z)
 \to\widehat B_4$ sending $(X,Y,z)$ to $(V,d^2,a^{-2})$.
Since $f$ has zero exponent sum in $X$, its evaluation at
$(Xz,Y)$ in this direct product is $(f(X,Y),1)$.
The displayed map therefore proves
\[
                      f(y_3,d^2)=f(V,d^2).              \tag{ABR.11}
\]
Antisymmetry gives $f(V,d^2)^{-1}=f(d^2,V)$ and
$f(a^2,b^2)^{-1}=f(b^2,a^2)$, converting (ABR.9)--(ABR.10)
to exactly (ABR.3).

\par\medskip\noindent\textbf{The pure gauge, quotient maps, and finite data.}\quad
The braid $V$ is pure, and the permutation of $a$ has order two.
Thus $a^{8c}$ and $V^{-4c}$ lie in the closed pure subgroup for
every profinite $c$; $g$ induces no branch-label permutation.
Capping the two boundary components kills their boundary twists,
in particular their product $\omega$.  Both the quilt action
and the braid action carry this central product to its
$\lambda$th power.  Hence the comparison descends to
$B_4/\langle\omega\rangle$ and to its actual pure subgroup
$P_4/\langle\omega\rangle=\Gamma_{0,5}$.
Its inner comparison there is precisely the image of $g$.

For a finite quotient $q$, the two local move evaluations are
\[
\begin{array}{c|c|c}
\text{move}&\widehat F_2\longrightarrow Q&\text{additional powers}\\ \hline
 a\to b&(x,y)\longmapsto(q(b)^2,q(a)^2)&q(a)^{-8c},\ q(b)^{8c},\ q(V)^\nu\\
 e\to d&(x,y)\longmapsto(q(d),q(V)^2)&1.
\end{array}                                                    \tag{ABR.12}
\]
The third evaluation needed to compare the actions is
$(x,y)\mapsto(q(d)^2,q(V))$.  The actual homomorphism
(ABR.6) followed by $q$ supplies (ABR.4); the rest consists only
of the discrete identities (ABR.5) and the two duality evaluations.
The direct-product map in (ABR.11) remains an actual map after
composition with $q$.  Thus every finite argument and its source
are specified.  The coefficient $c$ is needed only modulo the
exponent of $Q$, and $\lambda$ only modulo twice that exponent
when the half-power $\nu$ is retained in the edge weight.
The coefficient is exactly the second coordinate of $f$ in the
index-two Schreier abelianization
$\langle x^2,y,xyx^{-1}\rangle^{\rm ab}$, reduced modulo that
exponent; this is a finite map with quotient size at most
$2\exp(Q)^3$.

The executed independent script
\path{work/tower_code/check_finite_braid_bridge.py}
uses the faithful Artin action on the free group of rank four.
It verifies all twelve discrete identities used above and
records both sides as literal braid words, including
$V=a^2y_3$, the two-chain relation, $(dV)^2=\omega a^2$,
and every commutation in (ABR.5).  The f-dependent equalities
(ABR.7)--(ABR.11) are proved symbolically above; the script is not
used to replace them by sample values.
\end{proof}
